\documentclass[11pt]{article}
\usepackage{amsmath,amssymb,amsthm,mathtools}
\usepackage[margin=1in]{geometry}
\usepackage{enumitem}
\usepackage{booktabs}
\usepackage{hyperref}

\newtheorem{theorem}{Theorem}
\newtheorem{lemma}{Lemma}
\newtheorem{definition}{Definition}
\newtheorem{proposition}{Proposition}
\newtheorem{warning}{Warning}
\newtheorem{obligation}{Obligation}

\newcommand{\AZ}{\mathsf A_Z}
\newcommand{\Kcmp}{\mathsf K_{\rm cmp}}
\newcommand{\Yfull}{Y^-_{\rm full}}
\newcommand{\Ycore}{Y^-_{\rm core}}
\newcommand{\Yfin}{Y^-_{N,T}}
\newcommand{\Ypub}{Y^-_{\rm pub}}
\newcommand{\Fix}{\operatorname{Fix}}
\newcommand{\Ran}{\operatorname{Ran}}
\newcommand{\Ker}{\operatorname{Ker}}
\newcommand{\tr}{\operatorname{tr}}

\title{Step 74: Weil Test Space and Anti-Invariant Response Space for the RH Membrane Route}
\author{Working note}
\date{}

\begin{document}
\maketitle

\section{Purpose}
The previous steps reduced the RH-facing membrane route to the chain
\[
  \AZ \preceq \Kcmp \preceq \Theta^-,
  \qquad \tr \Theta^- \to 0,
\]
plus a separating anti-invariant zero readout and a fixed/exhaustive zero ledger.
The load-bearing bridge is the Douglas gate
\[
  V_Z = T W_{\rm cmp},\qquad \|T\|\le 1,
\]
where \(\AZ=V_Z^*V_Z\) and \(\Kcmp=W_{\rm cmp}^*W_{\rm cmp}\).

This note identifies the spaces on which such a statement must live.  It does \emph{not} construct the zeta carrier and does not prove RH.  Its goal is to prevent four overreads:
\begin{enumerate}[label=(\roman*)]
\item positivity on a finite test family is not positivity on the completed response space;
\item trace/scalar shadows are not Loewner domination;
\item a dense test class promotes only through closability/core records;
\item a signed explicit formula is not yet a positive feature-map carrier.
\end{enumerate}

\section{The four response levels}
Let \(J(s)=1-\overline{s}\) be the completed anti-linear symmetry.  In centered coordinates \(s=1/2+u\), it acts by \(u\mapsto -\overline u\).  The anti-invariant direction detects displacement from the fixed locus.

The RH route must distinguish four response spaces.

\begin{definition}[Full completed anti-invariant response space]
The full response space \(\Yfull\) is the Hilbert completion of the anti-invariant response domain on which the completed carrier form \(q_{\rm cmp}\) and the zero displacement form \(q_Z\) are both defined as closed or closable quadratic forms.  It is the only space on which a full Douglas bridge may be accepted.
\end{definition}

\begin{definition}[Dense Weil core]
A Weil core \(\Ycore\subset \Yfull\) is a dense subspace generated by admissible test functions for the completed explicit formula.  Depending on the carrier, this may be represented by compactly supported smooth log-side functions, Paley--Wiener spectral test functions, de Branges kernel vectors, or adelic test vectors.  It is accepted as load-bearing only if it is a form core for the completed feature form.
\end{definition}

\begin{definition}[Finite/truncated response family]
A finite response family \(\Yfin\subset\Ycore\) is a finite-dimensional span of selected tests, truncated prime windows, or truncated zero windows.  Positivity or domination on \(\Yfin\) is support-only unless it comes with an exhaustive tail bridge.
\end{definition}

\begin{definition}[Public trace shadow]
A public response shadow \(\Ypub\) is a scalar, trace, moment, or finite public readout of the response.  It may descend from a lawful witness by congruence, but it does not promote back to a carrier-level witness without reconstruction and defect accounting.
\end{definition}

\section{Form version of the Douglas bridge}
Let \(\mathcal C\) be a complex vector space, intended to be \(\Ycore\).  Let
\[
  q_Z[y]=\|V_Zy\|^2,
  \qquad
  q_{\rm cmp}[y]=\|W_{\rm cmp}y\|^2
\]
be nonnegative quadratic forms on \(\mathcal C\).  The desired inequality is
\[
  q_Z[y]\le q_{\rm cmp}[y]
  \qquad (y\in \mathcal C),
\]
with enough closability to extend it to the completed response space.

\begin{theorem}[Core domination promotes to full Douglas domination]
Let \(q_{\rm cmp}\) be a closable nonnegative form on \(\mathcal C\), and let \(\Yfull\) be the Hilbert space obtained by completing \(\mathcal C/\Ker q_{\rm cmp}\) in the norm \(q_{\rm cmp}^{1/2}\).  Suppose that \(q_Z\) is a nonnegative form satisfying
\[
  q_Z[y]\le q_{\rm cmp}[y]
  \qquad (y\in \mathcal C).
\]
Then \(q_Z\) descends to a bounded quadratic form on \(\Yfull\) with norm at most one.  Equivalently, there is a contraction \(T\) such that
\[
  V_Z = T W_{\rm cmp}
\]
on the completed form domain.  In operator notation,
\[
  \AZ\preceq \Kcmp .
\]
\end{theorem}

\begin{proof}
If \(q_{\rm cmp}[y]=0\), then \(q_Z[y]=0\), so \(q_Z\) descends to the quotient \(\mathcal C/\Ker q_{\rm cmp}\).  The inequality says the descended quadratic form is bounded by the identity with respect to the \(q_{\rm cmp}\)-norm.  Therefore it extends continuously to the completion.  In feature-map language, define \(T(W_{\rm cmp}y)=V_Zy\) on \(\Ran W_{\rm cmp}\).  The inequality makes this well-defined and contractive; extend by continuity to the closure.  This is the Douglas factorization in form language.
\end{proof}

\begin{warning}[Core positivity is not finite positivity]
If domination is checked only on \(\Yfin\), then the theorem does not apply.  The missing record is a proof that \(\Yfin\) is an exhaustive core, or a finite-window tail estimate that promotes \(\Yfin\) to the completed response space.
\end{warning}

\begin{warning}[Null form legality]
The bridge fails if \(q_{\rm cmp}\) has a null direction seen by \(q_Z\).  In the framework this is the zero-side version of null-mode legality: no anti-invariant zero displacement may live in a carrier-zero direction unless it is explicitly quotiented or annihilated.
\end{warning}

\section{Candidate Weil core records}
This note does not choose the zeta carrier.  It records three candidate ways to represent the core.

\subsection{Log-side core}
Work in logarithmic coordinate \(t=\log x\).  A provisional core is
\[
  \mathcal C_{\log}=C_c^\infty(\mathbb R)
\]
or a symmetry-stable subspace adapted to the completed involution.  Shift features have the form
\[
  (W_\mu f)(a,t)=\sqrt{\mu(a)}(f(t+a)-f(t)),
\]
with positive symbol
\[
  \Psi_\mu(\xi)=\int \mu(a)|e^{ia\xi}-1|^2\,da\ge0.
\]
This representation is convenient for prime-power and archimedean shift candidates, but requires support/tail and archimedean renormalization records.

\subsection{Spectral Paley--Wiener core}
A second core is a Paley--Wiener class of spectral test functions on the completed spectral variable.  This is closer to classical explicit formula presentations.  Its burden is to identify the anti-invariant response sector and to prove that the completed signed formula equals \(q_{\rm cmp}-q_Z\) on a form core.

\subsection{Carrier-native cores}
A third core may arise from a specific carrier: de Branges spaces, adelic test spaces, or a Selberg/trace-formula carrier.  Such a core is acceptable only after the exact package, null quotient, and all-six records are declared.

\section{The exact response-space obligation}
The RH O1 obligation can now be stated in layers.

\begin{obligation}[O1a: response space]
Construct \(\Yfull\), \(\Ycore\), and the anti-invariant projection \(P_-\), with \(\Ycore\) a core for the completed carrier form.
\end{obligation}

\begin{obligation}[O1b: zero displacement form]
Construct \(V_Z:\Ycore\to\mathcal Z\) so that
\[
  q_Z[y]=\|V_Zy\|^2
\]
is the anti-invariant zero-displacement form.  The readout must separate the fixed locus for the completed zero ledger.
\end{obligation}

\begin{obligation}[O1c: completed feature form]
Construct
\[
  W_{\rm cmp}=\begin{bmatrix}W_{\rm pr}\\ W_\infty\\ W_{\rm pole}\\ W_{\rm tail}\end{bmatrix}
\]
with positive completed feature form
\[
  q_{\rm cmp}[y]=\|W_{\rm cmp}y\|^2.
\]
Signed prime/gamma/pole formulas must be grouped into a positive representation before they can be used here.
\end{obligation}

\begin{obligation}[O1d: form-core domination]
Prove
\[
  q_Z[y]\le q_{\rm cmp}[y]
  \qquad (y\in\Ycore),
\]
with null-mode legality.  Then the core theorem gives the Douglas bridge on \(\Yfull\).
\end{obligation}

\section{Status classification}
\begin{center}
\begin{tabular}{lll}
\toprule
Record & What is proved & Status \\
\midrule
Full completed response & \(\AZ\preceq\Kcmp\) on \(\Yfull\) & accepted bridge \\
Dense core + closability & domination on form core extends & accepted after core record \\
Finite tests only & domination on \(\Yfin\) & support-only \\
Public trace shadow & scalar/trace balance & shadow-only \\
Signed explicit formula & identity of signed terms & incomplete until positive features \\
\bottomrule
\end{tabular}
\end{center}

\section{RH specialization}
For RH,
\[
  J(s)=1-\overline s,
  \qquad
  \Fix(J)=\{s:\Re(s)=1/2\}.
\]
The anti-invariant scalar readout
\[
  \psi_-(s)=\Re(s)-\frac12
\]
separates the fixed locus.  But the Douglas bridge is not about this scalar readout alone.  It requires a completed anti-invariant response space \(\Yfull\), a zero displacement form \(q_Z\), and a completed positive carrier feature form \(q_{\rm cmp}\), with
\[
  q_Z\le q_{\rm cmp}
\]
on the full response domain or on a closable form core.

\section{Nonclaims}
This note does not prove the Weil bridge.  It does not prove RH.  It does not claim that finite test positivity, trace equality, Li-type scalar positivity, or a signed explicit formula automatically imply the Douglas gate.  It only identifies the spaces and core-extension records needed before such a bridge can be honestly attempted.

\end{document}
